\titledquestion{Another definition of Brownian motion}
\begin{parts}
    % @QUESTION
    \part Show that a process 
    $(X_t)_{t\in \R^+}$ is a Brownian motion if and only if:
        \begin{itemize}
        \item[i')] $X_0=0$,
        \item[ii')] $(X_t)_{t\in \R^+}$ is a centered Gaussian process,
        \item[iii')] ${\rm Cov}(X_s,X_t)=\inf(s,t).$
        \end{itemize}
    % @QUESTION_END
    \begin{xsolution}
    % @SET_START
    Let us start with the direction: Brownian $\Rightarrow$ Gaussian process.
    % @LIST_START
    \begin{enumerate}
        \item Independence structure of increments :  
           % @ATOM
           % @PRECOND
           For \(s \leq t\), we decompose \(X_t = X_s + (X_t - X_s)\).
           % @ARGUMENT
           Independence of \(X_s\) and \(X_t - X_s\) 
           % @ARGUMENT_END
           implies that the covariance reduces to :  
           \[
           \mathrm{Cov}(X_s, X_t) = \mathrm{Cov}(X_s, X_s) + \mathrm{Cov}(X_s, X_t - X_s) = \mathrm{Var}(X_s) + 0 = s
           \]  
           which directly gives 
           % @OUTCOME
           \(\mathrm{Cov}(X_s, X_t) = \min(s,t)\).
           % @ATOM_END
        \item Gaussian character :  
           % @ATOM
           % @PRECOND
           Standard Brownian motion has Gaussian increments and 
           % @PRECOND
           \(X_0 = 0\), 
           % @PRECOND_END
           so any finite linear combination of \((X_{t_1}, \dots, X_{t_n})\) is 
           % @ARGUMENT
           a linear combination of independent Gaussian increments thus it is Gaussian.  
           % @OUTCOME
           The process is therefore centered Gaussian.
           % @ATOM_END
    \end{enumerate}    
    %These two points essentially exploit \emph{independence of increments} and \emph{normality} of increments.
    % @LIST_END
    Then the converse:
    % @LIST_START
    \begin{enumerate}
    \item Computation of covariance of increments :  
       % @ATOM
       % @PRECOND
       Using property ($iii'$)
       % @PRECOND
       Compute $\text{Cov}(X_t - X_s, X_v - X_u)$
       for \(s < t \leq u < v\).
       % @ARGUMENT:CALCUL
       \[
       \text{Cov}(X_t - X_s, X_v - X_u) = \min(t,v) - \min(t,u) - \min(s,v) + \min(s,u)
       \]
       This always gives 0, proving that 
       % @OUTCOME
       disjoint increments are uncorrelated, hence independent 
       % @OUTCOME_END
       (since Gaussian vector).
       % @ATOM_END
    \item Computation of variance of increments :  
       % @ATOM
       % @PRECOND
       Compute the variance of an increment of $X$:
       % @PRECOND_END
       % This is a calculus step
       For \(s < t\) :
       \[
       \text{Var}(X_t - X_s) = t + s - 2s = t - s
       \]
       % but the argument is
       % @ARGUMENT
       thus the law depends only on \(t-s\).
       % @OUTCOME
       This shows that the increments are stationary 
       % @OUTCOME
       and have law \(\mathcal{N}(0, t-s)\).
       % @ATOM_END
    These two points:
    % @SET_START
    % @ATOM
    % @PRECOND
    disjoint increments are independent, and
    % @PRECOND
    the law of increments is Gaussian centered of variance $t-s$, i.e. \(\mathcal{N}(0, t-s)\),
    % @ARGUMENT
    plus the assumption of continuity of paths
    % @ARGUMENT_END
    give 
    % @OUTCOME 
    the classical definition of Brownian motion.
    % @ATOM
    % @PRECOND
    disjoint increments are independent, and
    % @PRECOND
    the law of increments is Gaussian centered of variance $t-s$, i.e. \(\mathcal{N}(0, t-s)\),
    % @ARGUMENT
    by taking the continuous version via Kolmogorov's theorem
    % @ARGUMENT_END
    give 
    % @OUTCOME 
    the classical definition of Brownian motion.
    % @SET_END
    \end{enumerate}
    % @LIST_END
    % @SET_END
    \end{xsolution}

    % @QUESTION
    \part Show that if $W_t$ is a Brownian motion, $ \{X_t=t\;W_{\frac
    1t}, t>0, X_0=0 \}$
    is a Brownian motion too.
    \begin{xsolution}
    % @SET_START
    % @ATOM
    For the continuity at 0,
    % @PRECOND
    compute $\lim_{t \to 0^+} t W_{1/t}$:
    % @PRECOND_END
    \[
    \lim_{t \to 0^+} t W_{1/t} = \lim_{s \to \infty} \frac{W_s}{s} = 0 \quad \text{a.s.}
    \]
    % @ARGUMENT
    Hence by the law of large numbers for Brownian motion:
    % @OUTCOME 
    $\lim_{t \to 0^+} t W_{1/t} = 0$ a.s.
    % @ATOM
    \\
    For the centered Gaussian character :  
    % @PRECOND
    For any arbitrary \(t_1, \dots, t_n\), the vector \((X_{t_1}, \dots, X_{t_n})\) 
    % @ARGUMENT
    is a linear combination of the \(W_{1/t_i}\) 
    % @OUTCOME 
    thus it is centered Gaussian.
    \\
    % @ATOM
    For the Brownian covariance :  
    % @ARGUMENT:CALCUL
    \[
    \mathbb{E}[X_t X_s] = ts \cdot \min\left(\frac{1}{t}, \frac{1}{s}\right) = \min(t,s),
    \]
    % @OUTCOME 
    i.e. $\mathbb{E}[X_t X_s] = \min(t,s)$.
    \\
    % @ATOM
    Finally
    % @PRECOND
    continuity: $\lim_{t \to 0^+} X_t = 0$ a.s.,
    % @PRECOND
    Gaussianity: the vector \((X_{t_1}, \dots, X_{t_n})\) is centered Gaussian for any set of times, and
    % @PRECOND
    covariance: $\mathbb{E}[X_t X_s] = \min(t,s)$,
    % @PRECOND
    adding \(X_0 = 0\)
    % @ARGUMENT
    we get by definition of a Brownian motion
    % @OUTCOME
    that \(X\) is a Brownian motion.
    % @SET_END
    \end{xsolution}
\end{parts}